% Part: many-valued-logic
% Chapter: three-valued-logics
% Section: introduction

\documentclass[../../../include/open-logic-section]{subfiles}

\begin{document}

\olfileid[fr]{mvl}{thr}{int}

\olsection{Introduction}

Si nous ajoutons la valeur~$\Undef$ à
$\True$ et à~$\False$, nous obtenons
une logique à trois valeurs. Bien qu'il
n'y ait qu'une valeur de vérité
supplémentaire, les façons de définir
les fonctions de vérité de $\lnot$,
$\land$, $\lor$ et $\lif$ sont
nombreuses. Une logique peut alors
retenir n'importe quelle combinaison
de ces fonctions. On peut, par exemple, désigner $\True$ seul
ou désigner à la fois $\True$ et~$\Undef$.

Nous présentons ici une sélection des
logiques à trois valeurs les plus
connues, les motivations qui les
sous-tendent et certaines de leurs
propriétés.

\end{document}
